\chapter{数列与级数}\label{ch:b2:series}

本科一年级分册中的数值级数理论在此沿三个方向深化：
其一是取值于巴拿赫空间的级数，绝对收敛在其中起关键作用；
其二是实级数更精细的判别法（阿贝尔求和）；
其三是\emph{可和族}——摆脱项的次序的求和——
以及二重和的富比尼定理与柯西乘积。
这些工具为其后各章的函数项级数
提供基础。

\section{赋范空间中的级数}

\begin{definition}\label{def:b2:series:def}
设 $(u_n)$ 为赋范空间 $E$ 中的数列，级数 $\sum u_n$
在其部分和收敛时收敛；若级数
\emph{绝对收敛}\index{绝对收敛!在巴拿赫空间中}，即
$\sum \norm{u_n} < \infty$。在\emph{巴拿赫}空间中，绝对收敛
蕴含收敛（见\cref{thm:b2:nvs:absoluteconvergence}）；
在非完备空间中，这一结论可能失效
（见\cref{exo:b2:series:9}）。
\end{definition}

\begin{example}\label{ex:b2:series:neumann}
在 $\mathcal{M}_n(K)$（或 $\mathcal{L}_c(E)$，$E$ 为巴拿赫空间）中：
对 $\vertiii A < 1$，\emph{诺伊曼级数} $\sum A^k$
绝对收敛于 $(I - A)^{-1}$（证明见\cref{exo:b2:nvs:5}）；
$\sum \frac{A^k}{k!}$ 绝对收敛于 $\eu^A$，对每个 $A$
均成立（见\cref{ex:b2:nvs:matrixexp}）。算子值的几何级数与指数级数，
其行为与数量情形的模型相同——
这正是巴拿赫框架的要义所在。
\end{example}

\section{阿贝尔求和}

\begin{theorem}[阿贝尔求和与判别法]\label{thm:b2:series:abel}
（分部求和）对数 $a_n$ 与向量 $b_n$，记 $B_n
= \sum_{k=0}^{n} b_k$：\index{阿贝尔求和}
\[
\sum_{n=0}^{N} a_n b_n
= a_N B_N - \sum_{n=0}^{N-1} (a_{n+1} - a_n) B_n .
\]
（阿贝尔判别法）若 $(a_n)$ 为实数列，单调递减趋于 $0$，且
其部分和 $B_n$（在巴拿赫空间中）\emph{有界}，则
$\sum a_n b_n$ 收敛。
\end{theorem}

\begin{proof}
恒等式的证明分步进行：取 $B_{-1} = 0$，写 $b_n = B_n -
B_{n-1}$，并作拆分：
\[
\sum_{n=0}^{N} a_nb_n
= \sum_{n=0}^{N} a_nB_n - \sum_{n=0}^{N} a_nB_{n-1}
= \sum_{n=0}^{N} a_nB_n - \sum_{n=0}^{N-1} a_{n+1}B_{n} ,
\]
对第二个和作指标替换 $n \mapsto n + 1$（$B_{-1}$
项为零）；合并公共求和范围 $0 \leq n \leq N-1$
得 $a_NB_N$ 加上 $\sum_{n\leq N-1}(a_n - a_{n+1})B_n$，
即所述公式。这是离散形式的分部积分，其中
以 $(B_n)$ 作为 $(b_n)$ 的原函数，以差 $a_{n+1}
- a_n$ 作为 $(a_n)$ 的导数。对判别法，设 $\norm{B_n}
\leq M$，则
边界项 $a_N B_N \to 0$；级数 $\sum (a_n -
a_{n+1})B_n$ 绝对收敛，因为
\[
\sum_n \norm{(a_{n+1} - a_n)B_n} \leq M \sum_n (a_n - a_{n+1})
= M a_0 < \infty
\]
（裂项相消，$a_n \downarrow 0$）。恒等式的两部分
都收敛，从而 $\sum a_n b_n$ 收敛。
\end{proof}

\begin{example}\label{ex:b2:series:sinn}
$\sum \frac{\sin n}{n}$ 收敛：取 $a_n = \frac1n \downarrow 0$，
且 $B_n = \sum_{k=1}^{n} \sin k$ 有界——事实上 $B_n =
\Im\sum_{k \leq n} \eu^{\iu k} = \Im\,\frac{\eu^{\iu}(\eu^{\iu n} -
1)}{\eu^{\iu} - 1}$，其模 $\leq \frac{2}{\abs{\eu^{\iu} - 1}}$。
它\emph{不}绝对收敛（$\abs{\sin n} \geq \sin^2 n
= \frac{1 - \cos 2n}{2}$，且 $\sum \frac{1 - \cos 2n}{2n}$ 发散，
因为 $\sum \frac{\cos 2n}{n}$ 由同一阿贝尔判别法收敛，而
$\sum \frac{1}{2n}$ 发散）。交错级数判别法是
$b_n = (-1)^n$ 时的特殊情形。
\end{example}

\begin{example}[收敛圆周上的阿贝尔判别法]\label{ex:b2:series:abelboundary}
对哪些复数 $z$ 满足 $\abs z = 1$，级数 $\sum_{n \geq 1}
\frac{z^n}{n}$ 收敛？当 $z = 1$ 时，它是调和级数，
发散。对圆周上的 $z \neq 1$，可用阿贝尔判别法，
取 $a_n = \frac1n \downarrow 0$，$b_n = z^n$，其部分和
与 $N$ 无关地有界：
\[
\Bigl|\sum_{n=1}^{N} z^n\Bigr|
= \Bigl|\frac{z(z^N - 1)}{z - 1}\Bigr|
\leq \frac{2}{\abs{z - 1}} .
\]
收敛——但从不绝对收敛（$\sum\frac1n$）。
同一级数在圆周上表现出多种行为：单点处发散，
其余各点处条件收敛。这正是幂级数的标准边界行为
（见\cref{ch:b2:powerseries}），此处仅以初等方法研究，
在 $z = -1$ 时它恰好给出交错
调和级数，而在 $z = \eu^{\iu\theta}$ 时，其实部与虚部
分别为级数 $\sum\frac{\cos
n\theta}{n}$ 与 $\sum\frac{\sin n\theta}{n}$，
即\cref{exo:b2:series:4}所讨论的级数。
\end{example}

\begin{example}[暗藏陷阱的交错级数]\label{ex:b2:series:trap}
级数 $\sum_{n\geq2} \dfrac{(-1)^n}{\sqrt n + (-1)^n}$ 是否收敛？
符号交替出现，项趋于 $0$——然而
交错判别法\emph{不}适用：其模
$\frac{1}{\sqrt n + (-1)^n}$ 并不单调递减（在每个奇数 $n$ 处
跳升）。改为展开：
\[
\frac{(-1)^n}{\sqrt n + (-1)^n}
= \frac{(-1)^n}{\sqrt n}\cdot
\frac{1}{1 + \frac{(-1)^n}{\sqrt n}}
= \frac{(-1)^n}{\sqrt n} - \frac{1}{n}
+ O\Bigl(\frac{1}{n^{3/2}}\Bigr).
\]
第一项收敛（交错判别法，如实应用于
$\frac1{\sqrt n}\downarrow0$），第三项绝对收敛——
但中间一项是发散的调和级数，从而
整个和\emph{发散}至 $-\infty$。最后的要点是：当
单调性不成立时，应展开至每一项
要么绝对收敛，要么是可直接判别的典型情形；隐藏的
$-\frac1n$ 对符号计数法而言不可见。
\end{example}

\begin{remark}[常见陷阱]\label{rem:b2:series:pitfalls}
（i）“通项趋于 $0$”说明不了任何问题：调和级数
发散。（ii）交错判别法要求\emph{单调递减}的模——
\cref{ex:b2:series:trap}正是典型的
反例，而\cref{exo:b2:series:1}中的第三个级数即为
巩固练习。（iii）条件收敛级数不可重排
（\cref{ex:b2:series:rearrange}），其柯西
乘积可能发散：对 $\sum\frac{(-1)^n}{\sqrt{n+1}}$
自乘后，其对角项满足
\[
\abs{c_k} = \sum_{m=0}^{k}
\frac{1}{\sqrt{(m+1)(k-m+1)}}
\geq (k+1)\cdot\frac{2}{k+2} \longrightarrow 2 \neq 0
\]
（每个因子由算术–几何平均不等式至多为 $\frac{k+2}2$），故 $\sum c_k$
发散 --- 至少一个因子的绝对收敛
（\cref{exo:b2:series:8}）是不可或缺的。（iv）可和性
按定义，可和性关乎\emph{绝对}的界：
并不存在条件可和族这回事。
\end{remark}

\section{可和族}

\begin{definition}\label{def:b2:series:summable}
设 $I$ 为可数指标集。族 $(u_i)_{i \in I}$ 由
\emph{非负实数}组成，
当有限部分和有界时，称之为\emph{可和族}\index{可和族}；其和为
\[
\sum_{i \in I} u_i = \sup_{F \subseteq I \text{ 有限}} \sum_{i
\in F} u_i \in \intcc{0}{+\infty} .
\]
实数、复数（或巴拿赫空间中的向量）构成的族，
当 $(\norm{u_i})$ 可和时该族亦可和；其和则由拆分为
正、负部分（或实、虚部分）来定义——等价地，其和为
$\sum_{n} u_{\sigma(n)}$ 在
$\sigma$ 取遍 $I$ 的全部枚举时的公共值（见下文）。
\end{definition}

\begin{method}[选择判别法]\label{met:b2:series:choosing}
面对 $\sum u_n$，依次检验：（1）若 $u_n \not\to 0$，则发散，
即止。（2）若各项同号，则作比较：求与之等价的级数
（\cref{ch:b2:comparison}），并将其置于
黎曼--伯特兰图上定位。（3）若符号交错，
且模\emph{单调递减}，则用交错判别法；若模
不单调，则将通项展开，直至每一项
要么绝对收敛，要么为易于判别的典型情形（\cref{ex:b2:series:trap}）。（4）
若符号呈振荡而有规律（$\sin
n\theta$、$\eu^{\iu n\theta}$、矩阵幂），用阿贝尔判别法，
且部分和有界。（5）绝对收敛总值得先检验：
它更强、与次序无关，并可用于柯西
乘积与富比尼定理。
\end{method}

\begin{example}[按对角线计数判别可和性]\label{ex:b2:series:diagonalcount}
对哪些 $s > 0$，族
$\bigl((m + n)^{-s}\bigr)_{m, n \geq 1}$ 可和？将有限部分和
按对角线 $m + n = k$ 分组：第 $k$ 条对角线
含 $k - 1$ 对，每对贡献 $k^{-s}$，故有限和
恰以下式为界（并可穷尽之）
\[
\sum_{k \geq 2} \frac{k - 1}{k^{s}} ,
\]
此为与 $k^{1-s}$ 等价的正项级数，
故可和当且仅当 $s - 1 > 1$，即 $s > 2$。二维指标
吞下整整一个幂次：一片项构成的平面较直线
“多一维地发散”——指标集的计数几何，
而非单项的大小，决定可和性。
（同样的计数表明 $\bigl((m^2 + n^2)^{-1}\bigr)$ 不
可和：在对角线 $m + n = k$ 上，每项至少为
$k^{-2}$，而 $(k-1)\cdot k^{-2}$ 之和像调和级数
一样发散。）
\end{example}

\begin{theorem}[可和性与次序]\label{thm:b2:series:rearrangement}
\begin{enumerate}
  \item 对非负族，其和在任一
        枚举下不变：$\sum_{i} u_i = \sum_{n=0}^{\infty}
        u_{\sigma(n)}$，对每一双射 $\sigma \colon \N \to I$。
  \item 实或复级数 $\sum u_n$ 为\emph{可交换收敛}的
        （即其任一重排均收敛且和不变）当且仅当其
        绝对收敛。
\end{enumerate}
\end{theorem}

\begin{proof}
（1）每个部分和 $\sum_{n \leq N} u_{\sigma(n)}$ 都是该族的一个有限
部分和（故 $\leq$ 上确界）；每个有限 $F$ 均
包含于某个 $\{\sigma(0), \dots, \sigma(N)\}$ 中（故上确界
$\leq$ 级数的极限）。两界相合。

（2）若 $\sum\abs{u_n} < \infty$：对任一重排 $\sigma$ 及
$\varepsilon > 0$，取 $N$ 使 $\sum_{n > N}\abs{u_n} \leq
\varepsilon$；当 $\sigma$ 已穷尽
$\intint{0}{N}$ 的序号之后，重排后的部分和与原极限之差
至多为 $\varepsilon$：和相同。若 $\sum
\abs{u_n} = \infty$ 但 $\sum u_n$ 收敛（实情形；复情形
逐分量可得）：正部与负部皆
发散，且可经重排趋于任一预先给定的极限
--- 黎曼定理，在\cref{exo:b2:series:5}中完成 ---
故可交换收敛不成立。
\end{proof}

\begin{example}[一个被当场抓获的重排]\label{ex:b2:series:rearrange}
交错调和级数之和为
$\sum_{n\geq1}\frac{(-1)^{n-1}}{n} = \ln 2$（本科一年级
分册）。将其重排为“一正两负”：
\[
1 - \frac12 - \frac14 + \frac13 - \frac16 - \frac18 + \frac15 -
\cdots
\]
每三项为一组，
\[
\frac{1}{2k-1} - \frac{1}{4k-2} - \frac1{4k}
= \frac{1}{4k-2} - \frac{1}{4k}
= \frac12\Bigl(\frac{1}{2k-1} - \frac1{2k}\Bigr),
\]
于是重排后的级数收敛于 $\frac12\ln 2$ --- 恰为
原和的一半，而所用各项完全相同。非绝对收敛的级数
会“记住”各项的次序；可和族
恰恰是不记次序的那一类。
\end{example}

\begin{example}[加括号无妨，去括号则不然]\label{ex:b2:series:grouping}
对\emph{收敛}级数的相邻项加括号，
绝不改变和：加括号后的部分和构成原部分和的
子列。去括号这一逆运算是禁止的：
\[
(1 - 1) + (1 - 1) + (1 - 1) + \cdots = 0 + 0 + \cdots = 0,
\]
然而去括号后的 $1 - 1 + 1 - 1 + \cdots$ 发散（部分
和在 $1$ 与 $0$ 之间振荡）。去括号只在
附加补偿性假设下才合法，例如各项
趋于 $0$ 且分组长度有界：此时两个
相邻加括号部分和之间，原部分和的偏移至多为
有界个数的 $o(1)$ 项之和，于是收敛性可以逆向传回。
这恰是使
\cref{ex:b2:series:rearrange}中的分块计算成为证明，而非
障眼法的条件。
\end{example}

\begin{theorem}[族的富比尼定理；柯西乘积]\label{thm:b2:series:fubini}
设 $(u_{m,n})_{(m,n) \in \N^2}$ 为可和的二重族（即
$\sup_F \sum_F \abs{u_{m,n}} < \infty$）。则\index{级数的
富比尼定理}
\[
\sum_{(m,n)} u_{m,n}
= \sum_{m=0}^{\infty}\Bigl(\sum_{n=0}^{\infty} u_{m,n}\Bigr)
= \sum_{n=0}^{\infty}\Bigl(\sum_{m=0}^{\infty} u_{m,n}\Bigr)
= \sum_{k=0}^{\infty}\Bigl(\sum_{m+n=k} u_{m,n}\Bigr),
\]
内层级数均绝对收敛。特别地，若 $\sum
a_m$ 与 $\sum b_n$ 绝对收敛，则其\emph{柯西
乘积}\index{柯西乘积}绝对收敛，且
\[
\Bigl(\sum_m a_m\Bigr)\Bigl(\sum_n b_n\Bigr)
= \sum_{k=0}^{\infty} c_k,
\qquad
c_k = \sum_{m=0}^{k} a_m b_{k-m} .
\]
\end{theorem}

\begin{proof}
\emph{非负情形。}每种分组（按行、按列或
对角线）算得相同的上确界：任一有限数对集合
包含在某个有限行块中（将每个分组和由有限部分和从下方控制，
由总和从上方控制），部分和的单调
收敛性便完成其余部分 --- 具体地，对于行：
由 $\sum_{m \leq M}\sum_{n \leq N} u_{m,n} \leq S$ 可得，令 $N
\to \infty$ 再令 $M \to \infty$，得 $\sum_m \sum_n u_{m,n} \leq S$；
反之，每个有限集 $F$ 都位于这样的矩形之中，故 $S \leq
\sum_m\sum_n u_{m,n}$。对角线情形：两个界同样成立，只需以三角形
代替矩形。

\emph{一般情形。}将其分为正部与负部（实部与
虚部），每部分都是可和的非负族；四种
分组方式在每部分上结果一致，因而在差上也一致；内层级数的
绝对收敛性由非负情形得出，
对 $\abs{u_{m,n}}$ 应用即可。

\emph{柯西乘积。}族 $u_{m,n} = a_m b_n$ 是可和的：
$\abs{a_mb_n}$ 的有限部分和以
$\bigl(\sum\abs{a_m}\bigr)\bigl(\sum\abs{b_n}\bigr)$ 为界。按行求和得
$\bigl(\sum a_m\bigr)\bigl(\sum b_n\bigr)$；按对角线得 $\sum_k
c_k$。
\end{proof}

\begin{example}[指数恒等式的严格处理]\label{ex:b2:series:exp}
对 $a, b \in \C$（或可交换矩阵）有：
\[
\Bigl(\sum_m \frac{a^m}{m!}\Bigr)\Bigl(\sum_n
\frac{b^n}{n!}\Bigr)
= \sum_k \sum_{m+n=k} \frac{a^m b^n}{m!\,n!}
= \sum_k \frac{(a + b)^k}{k!},
\]
按对角线应用二项式定理：$\eu^a \eu^b = \eu^{a+b}$
--- 仅由级数本身导出的 $\exp$ 的函数方程。
（交换性只用于二项式步骤；对于不可交换的
矩阵，该恒等式确实不成立，
见\cref{ch:b2:diffeq}。）
\end{example}

\begin{example}[作为计算工具的柯西乘积]\label{ex:b2:series:squaresum}
由几何级数和\cref{exo:b2:series:2}中的 $\sum_{n
\geq 1} nz^n = \frac{z}{(1-z)^2}$（$\abs z < 1$），再作一次
柯西乘积即可完成二阶矩的计算。将 $\sum_m
mz^m$ 与 $\sum_n z^n$ 相乘：对角线上的系数为 $\sum_{m=0}^k
m = \frac{k(k+1)}2$，于是
\[
\frac{z}{(1-z)^3} = \sum_{k\geq0}\frac{k(k+1)}{2}\,z^k ,
\]
再由恒等式 $n^2 = 2\cdot\frac{n(n+1)}2 - n$ 可得
\[
\sum_{n\geq1} n^2z^n = \frac{2z}{(1-z)^3} - \frac{z}{(1-z)^2}
= \frac{z(1+z)}{(1-z)^3} .
\]
取 $z = \frac12$，得 $\sum_{n\geq1}\frac{n^2}{2^n} =
\frac{\frac12\cdot\frac32}{\frac18} = 6$ --- 闭式值，
全程无需求导，只是把绝对收敛级数
像多项式一样相乘。同样的恒等式拼接
可算出每个 $\sum n^dz^n$，概率论学者会认出
几何分布的二阶阶乘矩
（见\cref{ch:b2:genfun}）。
\end{example}

\begin{example}[一个二重求和的计算]\label{ex:b2:series:doublesum}
设实数 $s > 1$，令 $\zeta(s) = \sum_{n\geq1} n^{-s}$。用
二重求和计数因数 --- 族 $(m^{-s}n^{-s})$ 在
$(m,n) \in (\N^*)^2$ 上可和（收敛正项
级数之积）--- 并按乘积 $q = mn$ 分组：
\[
\zeta(s)^2 = \sum_{m,n} \frac{1}{(mn)^s}
= \sum_{q=1}^{\infty} \frac{d(q)}{q^s},
\]
其中 $d(q)$ 为 $q$ 的因数个数。可和族
使组合问题化为分析问题。
\end{example}

\begin{example}[用富比尼定理求值：$\sum_n (\zeta(n) - 1) = 1$]\label{ex:b2:series:zetasum}
对整数 $n \geq 2$，$\zeta(n) - 1 = \sum_{k \geq 2} k^{-n}$。
双重族 $(k^{-n})_{k, n \geq 2}$ 可和：
先对几何列求和，
\[
\sum_{k\geq2}\sum_{n\geq2} \frac{1}{k^n}
= \sum_{k\geq2} \frac{1/k^2}{1 - 1/k}
= \sum_{k\geq2} \frac{1}{k(k-1)} = 1
\]
（裂项相消），且各项皆正，故
由\cref{thm:b2:series:fubini}许可改按行求和：
\[
\sum_{n\geq2}\bigl(\zeta(n) - 1\bigr) = 1 .
\]
无穷多个 $\zeta$ 值，各看上去像超越数，
它们的尾项之和恰为 $1$。小结：当
二重和各项为正时，应选取使其能够化简的次序计算
——此处列为几何级数，行则难以处理，
而富比尼定理使化简结果可转移到另一次序。
\end{example}

\begin{example}[几何级数解方程]\label{ex:b2:series:neumannsolve}
在巴拿赫空间 $\bigl(C(\intcc01),
\norm\cdot_\infty\bigr)$ 中解方程 $x - K(x) = y$，其中 $K(f)$ 为
常值函数 $\frac12\int_0^1 f$。算子范数为
$\vertiii K \leq \frac12 < 1$，故可用诺伊曼级数
（见\cref{ex:b2:series:neumann}）：$x = \sum_{n\geq0} K^n(y)$。
计算迭代：$K(y) = \frac12\int_0^1 y$（为常数），
而将 $K$ 作用于常数 $c$ 得 $\frac c2$，因此
$K^n(y) = \frac{1}{2^{n-1}}\cdot\frac12\int_0^1 y$，$n \geq
1$。对几何常数求和：
\[
x = y + \Bigl(\int_0^1 y\Bigr)
\sum_{n\geq1}\frac{1}{2^n} = y + \int_0^1 y .
\]
验证：$x - K(x) = y + \int y - \frac12\bigl(\int y + \int
y\bigr) = y$。无穷级数，有限答案，以及
一行即可验证——几何级数是一种求逆算法，
而不仅是收敛性的陈述。
\end{example}

\begin{example}[利用部分分式裂项相消]\label{ex:b2:series:telescope}
精确求和并不常见，裂项相消是其主要来源。
分解：
\[
\frac{1}{n(n+1)(n+2)}
= \frac{1}{2}\Bigl(\frac{1}{n(n+1)} -
\frac{1}{(n+1)(n+2)}\Bigr),
\]
（通分即可验证），故
部分和相消：
\[
\sum_{n=1}^{N}\frac{1}{n(n+1)(n+2)}
= \frac12\Bigl(\frac{1}{1\cdot2} -
\frac{1}{(N+1)(N+2)}\Bigr)
\longrightarrow \frac14 .
\]
同一模式——将通项写成 $c(u_n - u_{n+1})$，
其中 $(u_n)$ 为显式数列，曾解出\cref{exo:b2:series:10}
（反正切题），并可求出每一个 $\sum\frac{1}{n(n+1)\cdots(n +
k)} = \frac{1}{k\cdot k!}$。在本阶段，若存在精确和，
则通项中通常藏着一个裂项相消。
\end{example}

\begin{remark}[本卷内的展望]\label{rem:b2:series:perspectives}
之后三章直接用到本章的内容。对于
\cref{ch:b2:funcseq}：$\sum f_n$ 的正规收敛即
$\sum\norm{f_n}_\infty$ 的绝对收敛，它是在巴拿赫空间
$\bigl(C, \norm\cdot_\infty\bigr)$ 中——本章的
\cref{thm:b2:nvs:absoluteconvergence} 换了身装束。对于
\cref{ch:b2:powerseries}：在收敛圆内
一切皆绝对收敛且可和，故柯西乘积与
重排均可自由进行（这正是幂级数能像多项式那样相乘的原因）；
在边界上则由阿贝尔判别法接手
（见\cref{ex:b2:series:abelboundary}）。对于\cref{ch:b2:genfun}：
概率母函数是幂级数，其
运算——独立变量之和对应乘积，
复合分布对应二重和——都
可由\cref{thm:b2:series:fubini}保证。可和族是
即将展开的分析的“法务部”。
\end{remark}

\begin{remark}[本章的应用]
凡含无穷和者皆经此处：幂级数
（\cref{ch:b2:powerseries}）实为乔装的可和族，
傅里叶系数经柯西乘积相乘，并
由帕塞瓦等式重排（见\cref{ch:b2:fourier}），概率
母函数（见\cref{ch:b2:genfun}）即为富比尼定理
用于期望的应用。本科三年级分册把可和族
纳入计数测度上的勒贝格积分之中——
在那里，\cref{thm:b2:series:fubini}成为
富比尼–托内利定理的特例。
\end{remark}

\section{练习}

\begin{exercise}[$\star$]\label{exo:b2:series:1}
判断下列级数的敛散性：$\sum \dfrac{\cos n}{n}$；$\;\sum \dfrac{(-1)^n}{\ln
n}$；$\;\sum \dfrac{(-1)^n}{n^{3/4} + \cos n}$ \emph{（参照本科一年级的陷阱：
交错判别法需要单调性）}。
\end{exercise}

\begin{exercise}[$\star$]\label{exo:b2:series:2}
证明：当 $\abs z < 1$ 时，$\sum_{n\geq1} n z^{n} =
\dfrac{z}{(1-z)^2}$，利用 $\sum z^n$ 与
其自身的柯西乘积。
\end{exercise}

\begin{exercise}[$\star\star$]\label{exo:b2:series:3}
（克罗内克型引理）设 $\sum b_n$ 为收敛的实级数。
证明：利用阿贝尔分部求和，$\dfrac{1}{n}\sum_{k=1}^{n} k\,b_k
\to 0$。
\end{exercise}

\begin{exercise}[$\star\star$]\label{exo:b2:series:4}
研究 $\sum \dfrac{\sin(n\theta)}{n^\alpha}$ 的收敛性
（$\theta \in \R$，$\alpha > 0$） --- 对于哪些 $(\theta, \alpha)$，
级数是绝对收敛、条件收敛，还是发散？
\end{exercise}

\begin{exercise}[$\star\star\star$]\label{exo:b2:series:5}
（黎曼重排）设 $\sum u_n$ 为收敛但非绝对收敛的实级数，
且 $\ell \in \R$。证明：
$\sum u_n$ 的某个重排收敛于 $\ell$。
\emph{（说明正项子级数与负项子级数均发散；
然后贪心地交替取项：先取正项，直至部分和超过
$\ell$，再取负项，直至部分和低于该数，如此反复；各项趋于
$0$，由此迫使级数收敛于 $\ell$。）}
\end{exercise}

\begin{exercise}[$\star\star$]\label{exo:b2:series:6}
证明下述族
$\Bigl(\dfrac{x^{m+n}}{m!\,n!}\Bigr)_{(m,n)\in\N^2}$ 对每个
$x \in \R$ 可和，并重新推导恒等式 $(\eu^x)^2 =
\eu^{2x}$，并沿对角线 $m + n =
k$ 对二重和分组。
\end{exercise}

\begin{exercise}[$\star\star$]\label{exo:b2:series:7}
证明族 $\bigl(\frac{1}{m^2 n^2}\bigr)_{m,n \geq 1}$
可和，并按 $\gcd$ 分组，记 $q = \gcd(m,n)$，
\[
\zeta(2)^2 = \sum_{q\geq1} \frac{1}{q^4}
\sum_{\substack{a,b \geq 1\\ \gcd(a,b)=1}} \frac{1}{a^2b^2}
= \zeta(4) \cdot S,
\]
其中 $S = \sum_{\gcd(a,b)=1} \frac{1}{a^2b^2}$，由此推出 $S =
\zeta(2)^2/\zeta(4)$。\emph{（每个数对 $(m,n)$ 唯一地写成
$(qa, qb)$，其中 $\gcd(a,b) = 1$。）}
\end{exercise}

\begin{exercise}[$\star\star\star$]\label{exo:b2:series:8}
（阿贝尔乘积定理，简化版）设 $\sum a_n$
绝对收敛，$\sum b_n$ 收敛。证明：它们的
柯西乘积 $\sum c_n$ 收敛，且 $\sum c_n = (\sum
a_n)(\sum b_n)$。\emph{（记 $C_N = \sum_{k\leq N} c_k = \sum_n
a_n B_{N-n}$，其中 $B$ 为 $b$ 的部分和；按
$n \leq N/2$ 是否成立分两种情形，利用 $(B_m)$ 的有界性及
$(a_n)$ 的绝对收敛尾部。）}
\end{exercise}

\begin{exercise}[$\star\star\star$]\label{exo:b2:series:9}
在（非完备的）空间 $E$ 中，其元素为最终为零的实数列，
赋以上确界范数，试给出一个绝对收敛的级数，
但在 $E$ 中不收敛。\emph{（取 $u_n = 2^{-n} e_n$，其中 $(e_n)$ 为
标准序列。）}
\end{exercise}

\begin{exercise}[$\star\star$]\label{exo:b2:series:10}
验证恒等式 $\arctan(n+1) - \arctan(n) =
\arctan\dfrac{1}{n^2 + n + 1}$，并由此求出下列和的精确值：
\[
\sum_{n=1}^{\infty} \arctan\frac{1}{n^2 + n + 1} .
\]
\end{exercise}

\begin{exercise}[$\star\star$]\label{exo:b2:series:11}
判断下列级数的敛散性（可利用等价量）：
\[
\sum_n \bigl(\sqrt{n+1} - \sqrt n\bigr)^{\alpha}
\ (\alpha > 0), \qquad
\sum_n \Bigl(1 - \cos\frac1n\Bigr), \qquad
\sum_n \Bigl(\eu - \Bigl(1 + \frac1n\Bigr)^{\!n}\Bigr).
\]
\end{exercise}

\begin{exercise}[$\star\star\star$]\label{exo:b2:series:12}
设 $(a_n)$ 为正且\emph{递减}的数列，
且 $\sum a_n$ 收敛。证明 $n\,a_n \to 0$ \emph{（用一段尾部
控制 $n a_{2n}$）}。说明逆命题不成立，并说明单调性是
本质的，并给出具体反例。
\end{exercise}

\section{问题：欧拉的 $\zeta(2) = \pi^2/6$，借助柯西的
余切和}

欧拉最著名的恒等式 $1 + \frac14 + \frac19 + \cdots =
\frac{\pi^2}6$，其证明完全初等，归功于
柯西：由棣莫弗公式可得一多项式，其根为
数 $\cot^2\frac{k\pi}{2n+1}$，由韦达定理可精确求出这些根之和，
并利用夹逼式 $\cot^2\theta < \frac{1}{\theta^2} <
1 + \cot^2\theta$，将 $\sum\frac1{k^2}$ 的部分和夹在
两个显式有理界之间。下面完整写出证明，
并用同样方法推出 $\zeta(4) = \frac{\pi^4}{90}$，再
借助伯特兰级数刻画收敛与发散之间的整个
分界 --- 并证明该分界上根本不存在
最慢收敛的级数。

\begin{problem}[{周末问题 --- $\zeta(2) = \pi^2/6$ 与
伯特兰全景}]
\label{pb:b2:series:1}
全文设 $n \geq 1$，$\theta_k = \dfrac{k\pi}{2n+1}$，其中 $k
= 1, \dots, n$；注意 $0 < \theta_k < \frac\pi2$。

\medskip
\noindent\textbf{第一部分 --- 余切恒等式。}
\begin{enumerate}
  \item 证明棣莫弗公式 $(\cos\theta +
        \iu\sin\theta)^m = \cos m\theta + \iu\sin m\theta$（$m
        \in \N$），并由此推出，对 $m = 2n + 1$，
        \[
        \sin\bigl((2n{+}1)\theta\bigr) = \sum_{j=0}^{n}
        (-1)^j\binom{2n+1}{2j+1}\cos^{2(n-j)}\theta\,
        \sin^{2j+1}\theta .
        \]
  \item 由此可得，对 $\theta \in \intoo{0}{\frac\pi2}$，
        \[
        \sin\bigl((2n{+}1)\theta\bigr) =
        \sin^{2n+1}\theta\; P_n(\cot^2\theta),
        \qquad
        P_n(x) = \sum_{j=0}^{n}
        (-1)^j\binom{2n+1}{2j+1}x^{\,n-j},
        \]
        一个次数为 $n$ 的多项式，首项系数为 $2n
        + 1$。
  \item 证明 $x_k = \cot^2\theta_k$，$k = 1, \dots, n$ 是
        $n$ 个\emph{互异}的 $P_n$ 的根 --- 从而即为全部
        根。
  \item 由韦达定理，证明下列精确恒等式
        \[
        \sum_{k=1}^{n} \cot^2\frac{k\pi}{2n+1}
        = \frac{n(2n-1)}{3}.
        \]
  \item 由此亦可得 $\displaystyle\sum_{k=1}^{n}
        \frac{1}{\sin^2\theta_k} = \frac{2n(n+1)}{3}$.
  \item 证明夹逼式：$\cot^2\theta < \dfrac1{\theta^2} <
        \dfrac{1}{\sin^2\theta}$，对 $\theta \in
        \intoo{0}{\frac\pi2}$，\emph{（由 $\sin\theta < \theta <
        \tan\theta$）}。
\end{enumerate}

\medskip
\noindent\textbf{第二部分 --- 夹逼收尾：欧拉
定理。}
\begin{enumerate}[resume]
  \item 对 $k = 1, \dots, n$ 求和第 6 题，取 $\theta =
        \theta_k$，证明：
        \[
        \frac{n(2n-1)}{3} \;<\;
        \frac{(2n+1)^2}{\pi^2}\sum_{k=1}^{n}\frac1{k^2}
        \;<\; \frac{2n(n+1)}{3}.
        \]
  \item 由此得出（\emph{欧拉定理，柯西证法}）：
        \[
        \zeta(2) = \sum_{k=1}^{\infty}\frac{1}{k^2} =
        \frac{\pi^2}{6}.
        \]
  \item 由夹逼式推出收敛速度，证明：
        \[
        \Bigl|\sum_{k=1}^{n}\frac1{k^2} - \frac{\pi^2}6\Bigr|
        = O\Bigl(\frac1n\Bigr),
        \]
        与 $\sum_{k>n}k^{-2} =
        \frac1n - \frac{1}{2n^2} + O(n^{-3})$ 的级数余项相符，
        见\cref{exo:b2:comparison:11}。
  \item 把方法再推进一层：利用
        $P_n$ 的第二个韦达函数，证明
        \[
        \sum_{k=1}^{n}\cot^4\theta_k =
        \Bigl(\frac{n(2n-1)}3\Bigr)^{\!2} -
        \frac{2n(2n-1)(2n-2)(2n-3)}{60}
        \;\sim\; \frac{8n^4}{45},
        \]
        并用 $\cot^4 < \theta^{-4} < (1 +
        \cot^2)^2$ 夹逼，得 $\zeta(4) = \dfrac{\pi^4}{90}$。
\end{enumerate}

\medskip
\noindent\textbf{第三部分 --- 收益。}
\begin{enumerate}[resume]
  \item 由 $\zeta(2) = \frac{\pi^2}6$ 推出：
        \[
        \sum_{k\geq0}\frac{1}{(2k+1)^2} = \frac{\pi^2}{8},
        \qquad
        \sum_{k\geq1}\frac{(-1)^{k-1}}{k^2} = \frac{\pi^2}{12}.
        \]
  \item 结合 \cref{exo:b2:series:7}，计算 $S =
        \sum_{\gcd(a,b)=1}\frac{1}{a^2b^2} =
        \frac{\zeta(2)^2}{\zeta(4)} = \frac52$，并解释
        $\frac{1}{\zeta(2)} = \frac{6}{\pi^2} \approx 0.608$ 为
        互素数对的密度（诚实地说明启发式，严格的
        计数是本科三年级分册的内容
        ）。
  \item （可验证的加速）由第 9 题的余项公式
        可得 $\sum_{k\leq n}k^{-2} + \frac1n - \frac1{2n^2} =
        \frac{\pi^2}6 + O(n^{-3})$。比较求得
        $\zeta(2)$ 的六位数字所需的工作量：直接求和与
        $n = 100$ 时的修正和（此时误差为
        $1.7\cdot10^{-7}$）。
  \item 手算第 4 题，取 $n = 1$ 与 $n = 2$（其值分别为
        $\cot^2\frac\pi3 = \frac13$ 与
        $\cot^2\frac\pi5 + \cot^2\frac{2\pi}5 = 2$），利用
        $\cos\frac\pi5 = \frac{1+\sqrt5}4$ 或数值
        计算。
  \item 证明配套恒等式
        \[
        \sum_{k=1}^{n}\tan^2\frac{k\pi}{2n+1} = n(2n+1)
        \]
        \emph{（数 $\tan^2\theta_k$ 是反向多项式
        $x^nP_n(1/x)$ 的根）}，并在 $n
        = 1$ 时验证。
\end{enumerate}

\medskip
\noindent\textbf{第四部分 --- 伯特兰概览。}设 $\alpha,
\beta \in \R$，考虑\emph{伯特兰级数}
\[
\sum_{n \geq 3} \frac{1}{n^{\alpha}(\ln n)^{\beta}} .
\]
\begin{enumerate}[resume]
  \item 证明：当 $\alpha > 1$ 时，级数收敛，
        无论 $\beta$ 取何值 \emph{（与
        $n^{-(1+\alpha)/2}$ 比较）}。
  \item 证明：当 $\alpha < 1$ 时级数发散，无论
        $\beta$ 取何值。
  \item 当 $\alpha = 1$ 时，利用级数--积分比较
        （\cref{thm:b2:comparison:seriesintegral}），取 $f(t) =
        \frac{1}{t(\ln t)^\beta}$，证明当且仅当 $\beta
        > 1$ 时收敛。
  \item 进一步迭代这一临界：证明 $\sum\frac{1}{n\ln n\,\ln\ln
        n}$ 发散，而 $\sum\frac{1}{n\ln n\,(\ln\ln n)^2}$
        收敛。
  \item 两个陷阱：判断下列级数的敛散性
        \[
        \sum_n \frac{1}{n^{1 + 1/\ln n}}
        \qquad\text{ 与 }\qquad
        \sum_n \frac{1}{n^{1 + 1/\ln\ln n}}
        \]
        \emph{（精确计算 $n^{1/\ln n}$；比较
        $n^{1/\ln\ln n}$ 与 $\ln n$ 的任意幂）}。
  \item （不存在最慢的收敛级数）设 $\sum a_n$ 为任一
        收敛级数，其中 $a_n > 0$，且 $R_n = \sum_{k
        \geq n}a_k$ 为其余项。证明 $\sum
        \frac{a_n}{\sqrt{R_n}}$ 仍 \emph{收敛}
        \emph{（与裂项和 $2(\sqrt{R_n} -
        \sqrt{R_{n+1}})$ 作比较）}，虽然 $\frac{a_n/\sqrt{R_n}}{a_n}
        \to \infty$，但每一收敛级数都被另一收敛级数
        严格控制。收敛的临界
        不是一条曲线，而是一团雾。
\end{enumerate}

\medskip
\noindent\textbf{第五部分 --- 交叉检验与综合。}
\begin{enumerate}[resume]
  \item （柯西凝聚）证明：若 $(a_n)$ 为正项递减数列，则
        $\sum a_n$ 收敛当且仅当 $\sum 2^k a_{2^k}$
        收敛。由此重新推出第 18 题的收敛临界。
  \item （慢收敛的代价）对 $\sum\frac1{n(\ln n)^2}$，
        用积分估计其余项，并证明：部分和累加至 $N = 10^6$ 时，
        误差仍大于 $0.07$：由理论确认的收敛，
        对数值计算可能毫无用处
        --- 与第 13 题形成对比。
  \item 判定下列级数的敛散性（每题用一句话说明理由）：
        $\sum\frac1{n\ln n}$, $\sum\frac1{n^{1.01}}$,
        $\sum\frac{(\ln n)^{100}}{n^{1.001}}$,
        $\sum\frac1{n(\ln n)(\ln\ln n)^{3}}$.
  \item （综合）每小题用一句话回答：棣莫弗如何把一个
        三角恒等式化为根之和可计算的多项式；夹逼准则在何处
        需要\emph{精确}的
        端点恒等式而非等价关系；哪一工具
        在\cref{ch:b2:comparison}中支撑了第四部分；以及
        第 21 题对“通用比较判别法”的梦想说了什么；并点出两座高峰：
        欧拉的
        $\zeta(2) = \frac{\pi^2}6$（及其上一层的
        $\zeta(4) = \frac{\pi^4}{90}$），以及伯特兰
        分类。指出 $\zeta(2)$ 将于何处再次证明：
        在\cref{ch:b2:fourier}中由帕塞瓦尔给出 --- 一个
        定理，两种文明。
\end{enumerate}
\end{problem}
